\subsection{现有定义的证明边界}
\subsection{三种证明模板}
存在性证明通常依次验证非空性、紧性与连续性；唯一性证明通常依次验证严格凸性或严格拟凹性；比较静态证明通常依次验证目标函数、约束和参数映射的正则性。把证明拆成模板有助于避免把不同层次的结论混在一起。

\subsection{边界条件的经济解释}
当价格趋近于零、收入趋近于零或配额趋近于上限时，内点条件可能失效。研究中应明确说明参数空间是否排除这些边界，以及结论能否连续延拓到边界。边界不是技术噪声，而是政策阈值和制度约束最可能出现的区域。

\subsection{小结}
本章的形式化主线可以压缩为：先规定可行性，再规定评价，再定义选择，最后检验预测。任何跳过其中一步的论证，都可能把制度约束、偏好判断与经验事实混为一谈。
